\section{Marco teórico}
\label{ch:marco_teorico}

El marco teórico que sigue se organiza en nueve ejes: integridad académica en evaluaciones digitales, inteligencia artificial generativa y detección textual, trazabilidad de autoría como enfoque conceptual, biometría conductual y autenticación continua, dinámica de tipeo, señales contextuales de interacción, detección de anomalías y revisión humana, construcción del conjunto de datos, y privacidad y protección de datos. Una síntesis final relaciona estos ejes con el problema definido en la Sección~\ref{ch:problema}.

\subsection{Integridad académica y evaluación digital}
\label{sec:integridad_academica}

La integridad académica funciona como condición de validez de la evaluación. Una instancia evaluativa debe permitir inferir, de manera razonable y justa, el aprendizaje alcanzado por el estudiante. El International Center for Academic Integrity la formula como un compromiso con seis valores fundamentales, lo que permite entenderla como una dimensión institucional y cultural que excede el control de infracciones~\parencite{ICAI2021}.

Las evaluaciones en línea presentan desafíos específicos. \textcite{Noorbehbahani2022} revisaron 58 publicaciones sobre deshonestidad en exámenes en línea publicadas entre 2010 y 2021 y organizaron la literatura en torno a tipos de fraude, motivaciones, métodos de detección y estrategias de prevención. \textcite{Newton2024} reportaron presencia significativa de conductas deshonestas en evaluaciones en línea, aunque destacaron limitaciones derivadas de la heterogeneidad metodológica de los estudios incluidos.

El \textit{online proctoring} surgió como respuesta tecnológica a estas dificultades. La revisión de \textcite{Nigam2021} describe sistemas de supervisión con componentes como autenticación, monitoreo de actividad, análisis de video, bloqueo de navegación y detección automatizada de eventos. Sin embargo, \textcite{Mutimukwe2023} advierten que estas tecnologías pueden involucrar recolección y circulación de datos personales o sensibles, por lo que requieren transparencia, minimización, seguridad, responsabilidad y criterios de equidad.

Estas discusiones muestran que la integridad académica en entornos digitales no puede abordarse únicamente mediante vigilancia o automatización. Requiere articulación entre diseño evaluativo, controles técnicos proporcionales, revisión humana y criterios institucionales claros.

\subsection{Inteligencia artificial generativa, autoría y detección textual}
\label{sec:ia_generativa}

La inteligencia artificial generativa introdujo nuevas formas de asistencia en tareas académicas, especialmente en actividades de producción escrita. Los modelos de lenguaje de gran escala que sustentan las herramientas actuales se basan en la arquitectura \textit{Transformer}, que reemplazó las redes recurrentes por mecanismos de atención para modelar secuencias de texto~\parencite{Vaswani2017}; al escalar esa arquitectura en cantidad de parámetros y de datos de entrenamiento, los modelos adquirieron la capacidad de resolver tareas nuevas a partir de unas pocas indicaciones, sin reentrenamiento específico~\parencite{Brown2020}. Esa capacidad permite producir respuestas académicas completas a partir de una consigna, lo que plantea el problema de autoría que se analiza a continuación. UNESCO sostiene que su incorporación en educación requiere marcos institucionales, alfabetización en IA, protección de datos y definiciones explícitas sobre usos permitidos, restringidos o no permitidos~\parencite{UNESCO2023}. El problema de la autoría no se reduce a identificar si un texto fue generado por IA, sino a establecer condiciones transparentes para evaluar el proceso y la responsabilidad del estudiante sobre la entrega.

Los detectores automáticos de texto generado por IA presentan limitaciones documentadas. \textcite{Liang2023} evaluaron detectores GPT y reportaron que clasificaban erróneamente textos de escritores no nativos de inglés como generados por IA, lo que plantea riesgos de equidad en contextos evaluativos. \textcite{Krishna2023} mostraron que la paráfrasis puede evadir detectores bajo condiciones experimentales. \textcite{Perkins2024} evaluaron seis detectores de GenAI y observaron reducciones de desempeño frente a técnicas simples de evasión.

Estas limitaciones no implican descartar todo análisis textual, sino evitar que una clasificación IA/no IA opere como evidencia única o concluyente. Una misma señal textual puede ser compatible con escritura humana asistida, corrección gramatical, traducción, paráfrasis, uso parcial de IA o generación completa. Por ello, la literatura propone desplazar el énfasis desde la detección aislada del producto final hacia señales complementarias del proceso de producción.

\subsection{Trazabilidad de autoría como enfoque conceptual}
\label{sec:trazabilidad}

En este trabajo, la trazabilidad de autoría se define operativamente como la posibilidad de registrar, organizar y analizar señales asociadas al proceso mediante el cual se produce una respuesta académica en un entorno digital. Esta definición se concentra en información observable dentro de la plataforma, especialmente señales conductuales de escritura y eventos contextuales de interacción.

La trazabilidad de autoría se diferencia de la detección de fraude. La detección de fraude supone establecer que existió una conducta académicamente indebida, lo cual depende de normas institucionales, reglas de la consigna, evidencia disponible y procedimientos de revisión. La trazabilidad, en cambio, documenta condiciones observables del proceso de producción de la respuesta. Esta distinción es consistente con las advertencias de la literatura sobre tecnologías de supervisión educativa, que recomiendan evitar decisiones automatizadas no revisables y sostener criterios de transparencia y proporcionalidad~\parencite{Mutimukwe2023}.

También se diferencia de la detección textual de IA. Mientras los detectores textuales se centran en el producto final, la trazabilidad de autoría incorpora señales del proceso. La literatura fundamenta esta orientación en las limitaciones reportadas en detectores de texto generado por IA y en la necesidad de complementar el análisis del contenido con evidencia procesual interpretable~\parencite{Liang2023,Krishna2023,Perkins2024}.

\subsection{Biometría conductual y autenticación continua}
\label{sec:biometria_conductual}

El reconocimiento biométrico puede definirse como la verificación o identificación automática de individuos a partir de características fisiológicas o conductuales. \textcite{Jain2004} proponen propiedades relevantes para evaluar modalidades biométricas: universalidad, unicidad, permanencia, medibilidad, rendimiento, aceptabilidad y resistencia a la elusión.

La biometría conductual se basa en patrones de comportamiento observables, como la forma de escribir, interactuar con dispositivos o ejecutar acciones digitales. \textcite{Yampolskiy2008} la clasifican como un conjunto de técnicas orientadas a analizar patrones de conducta humana con fines de reconocimiento o verificación. En evaluaciones digitales, su interés radica en que ciertas señales pueden capturarse durante la interacción del estudiante con la plataforma, sin recurrir necesariamente a cámara o micrófono.

La autenticación continua extiende la lógica de verificación más allá del inicio de sesión. En lugar de validar al usuario únicamente al ingresar, compara señales observables durante la sesión con un perfil previamente construido~\parencite{Yampolskiy2008}. Para el campo educativo, este concepto permite fundamentar el análisis de consistencia conductual, que admite dos referencias: el historial del propio estudiante (línea base individual o intrausuario) o el comportamiento de un grupo comparable que rinde la misma evaluación (línea base poblacional o de cohorte).

Junto con el teclado, el dispositivo apuntador constituye otra fuente de señales conductuales. La dinámica del mouse (\textit{mouse dynamics}) caracteriza a un usuario a partir de la velocidad, la dirección, las pausas y los clics de sus movimientos, y fue propuesta como modalidad biométrica para la verificación de identidad~\parencite{Ahmed2007}. Trabajos posteriores la evaluaron para autenticación continua y señalaron como una de sus principales dificultades la variabilidad del comportamiento entre sesiones, lo que obliga a tratarla, igual que la dinámica de tipeo, como un indicio probabilístico~\parencite{Shen2013}.

Las modalidades conductuales presentan variabilidad. La literatura sobre \textit{keystroke dynamics} advierte que el comportamiento puede verse afectado por el tipo de tarea, el dispositivo, el contexto y el estado del usuario~\parencite{Teh2013}. Por ello, las señales biométricas conductuales se interpretan como indicadores probabilísticos y no como prueba concluyente de identidad, autoría o infracción.

\subsection{Dinámica de tipeo en respuestas escritas}
\label{sec:dinamica_tipeo}

La dinámica de tipeo o \textit{keystroke dynamics} estudia los patrones temporales de escritura producidos al utilizar un teclado. Entre las variables habituales se encuentran el tiempo de permanencia de una tecla (\textit{dwell time}), el tiempo de transición entre teclas (\textit{flight time}), las latencias entre pulsaciones, el ritmo de escritura, las pausas y los patrones de corrección. La revisión de \textcite{Teh2013} sistematiza esta modalidad biométrica y describe su uso para autenticación, con métricas como FAR (\textit{False Acceptance Rate}), FRR (\textit{False Rejection Rate}) y EER (\textit{Equal Error Rate}).

La literatura distingue escenarios de texto fijo y texto libre. \textcite{Killourhy2009} compararon catorce algoritmos de detección de anomalías en un \textit{benchmark} de contraseña fija con 51 participantes. Para texto libre, \textcite{Bergadano2002} propusieron un método de autenticación basado en patrones de escritura no restringidos. Esta distinción es relevante porque las respuestas abiertas de una evaluación se aproximan más al texto libre, escenario de mayor variabilidad metodológica.

En años recientes se exploraron modelos de mayor escala. \textcite{Acien2022} presentaron TypeNet, un modelo de \textit{deep learning} para biometría de tipeo evaluado en bases amplias, con distintos niveles de error según el tipo de teclado y el escenario de captura. En el plano educativo, Kundu et al. y Mehta et al. estudiaron la posibilidad de usar dinámica de tipeo para diferenciar escritura genuina de escritura asistida por modelos de lenguaje de gran escala (LLM) en condiciones experimentales, una línea de investigación emergente~\parencite{Kundu2024, Mehta2025}.

En evaluaciones digitales, la dinámica de tipeo permite fundamentar una noción de consistencia conductual: comparar una sesión con una línea base y estimar desviaciones. Las dos referencias posibles tienen requisitos distintos: la línea base individual exige sesiones previas de la misma persona, mientras que la poblacional solo requiere un conjunto de sesiones comparables, a cambio de no registrar las variaciones respecto del patrón habitual de cada estudiante. En ambos casos, la desviación estimada constituye un indicio técnico complementario, sujeto a revisión contextual.

\subsection{Señales contextuales de interacción}
\label{sec:senales_contextuales}

Además de las señales biométricas, una plataforma de evaluación puede registrar eventos contextuales producidos dentro de la interfaz, como pérdida de foco de la ventana, inactividad prolongada, pegado de texto o cambios bruscos en el ritmo de edición. Estas señales no identifican a una persona y, por lo tanto, no deben confundirse con biometría conductual. Su función es describir condiciones observables del proceso de interacción.

La literatura sobre deshonestidad en evaluaciones en línea ha explorado la combinación de señales técnicas y conductuales para detectar comportamientos anómalos. Tiong, Lee y Lim propusieron un enfoque que combina información de protocolo de Internet y clasificación conductual para la detección de conductas indebidas en evaluaciones en línea, luego publicado en una versión revisada por pares~\parencite{Tiong2022, Tiong2024}. Este tipo de trabajos respalda la pertinencia de analizar señales de interacción como insumo de una interpretación contextual.

Por razones metodológicas, las señales contextuales se describen en términos de lo observado. Un pegado extenso indica que un fragmento de texto fue incorporado mediante portapapeles; una pérdida de foco indica que la ventana de evaluación dejó de estar activa; una pausa prolongada indica ausencia de interacción registrada durante un intervalo. Ninguna de estas señales equivale necesariamente a copia, uso indebido de IA o falta de autoría: su valor depende de la consigna, de las reglas de la evaluación y de la revisión docente.

\subsection{Detección de anomalías y revisión humana}
\label{sec:deteccion_anomalias}

Los problemas de consistencia conductual pueden formularse como detección de anomalías cuando se dispone principalmente de ejemplos del comportamiento esperado y no de una representación completa de todas las formas posibles de asistencia externa, suplantación o comportamiento atípico. \textcite{Chandola2009} definen la detección de anomalías como la identificación de patrones que no se ajustan al comportamiento esperado de un conjunto de datos.

En biometría conductual aplicada a evaluaciones, la literatura recurre a esta formulación cuando no es posible entrenar un clasificador universal de copia/no-copia: los modelos de una sola clase y los detectores de anomalías son pertinentes cuando la clase negativa es amplia, cambiante o difícil de etiquetar de forma exhaustiva~\parencite{Chandola2009}.

Entre las técnicas no supervisadas, \textit{Isolation Forest} parte de una idea distinta de la de los métodos basados en distancia o densidad: en lugar de modelar lo normal, aísla cada observación mediante particiones aleatorias del espacio de variables, y como las observaciones anómalas son pocas y diferentes, quedan aisladas con menos particiones que las típicas~\parencite{Liu2008, Liu2012}. Las comparaciones empíricas entre algoritmos no supervisados muestran que el desempeño relativo depende del conjunto de datos y que ningún método resulta superior en todos los casos~\parencite{Goldstein2016, Campos2016}. Cuando se conocen las etiquetas de un conjunto de prueba, el desempeño de estos detectores suele resumirse con la curva ROC y el área bajo ella, que no dependen de un umbral particular~\parencite{Fawcett2006}, o con la curva de precisión y exhaustividad, más informativa cuando las anomalías son minoría~\parencite{Saito2015}.

La salida de estos modelos debe interpretarse como un indicador revisable. En contextos educativos, una decisión automatizada sin revisión humana puede ser problemática por el impacto potencial sobre estudiantes y por las limitaciones de las señales utilizadas. La literatura sobre \textit{proctoring} y privacidad recomienda transparencia, minimización, rendición de cuentas y criterios claros de uso cuando se recolectan y procesan datos de estudiantes~\parencite{Mutimukwe2023}.

La utilidad de estos indicadores para una decisión humana depende, además, de su explicabilidad, entendida como la posibilidad de presentar en términos comprensibles para una persona qué señales y qué métricas los originaron y con qué criterio~\parencite{DoshiVelez2017}. La literatura sobre explicación en inteligencia artificial señala que las personas valoran mejor un resultado cuando se les presentan sus causas concretas y se lo contrasta con alternativas, que cuando se les ofrece solo una probabilidad~\parencite{Miller2019}; una explicación de ese tipo permite juzgar la pertinencia de un indicador y considerar explicaciones alternativas antes de decidir. Desde esta perspectiva, un indicador que informa la regla, la métrica y el umbral que lo originaron ofrece una explicación causal del tipo que describe esa literatura, mientras que la salida de un modelo estadístico de anomalías no ofrece por sí misma esa explicación: puede acompañarse de información descriptiva sobre las variables en que la observación se aparta de la población, pero esa información no equivale a explicar la decisión del modelo.

\subsection{Construcción del conjunto de datos y factibilidad}
\label{sec:marco_dataset}

La detección de anomalías en datos de interacción requiere registros obtenidos en condiciones comparables a las de uso. Los conjuntos públicos de biometría de tipeo, como los utilizados para comparar modelos de autenticación~\parencite{Killourhy2009, Acien2022}, se obtuvieron en tareas de transcripción o de escritura libre ajenas a un examen, y no incluyen señales de navegación, cambios de foco ni eventos de pegado, ni etiquetas sobre la condición en que se produjo cada sesión. Por ese motivo, los trabajos sobre evaluaciones en línea suelen construir sus propios conjuntos de datos con participantes que resuelven tareas en condiciones controladas~\parencite{Kundu2024, Mehta2025}.

La literatura sobre evaluación de detectores de anomalías distingue tres fuentes de datos, de realismo creciente. Los datos sintéticos, generados a partir de parámetros declarados, permiten conocer de antemano qué observaciones son anómalas y estudiar el comportamiento de un algoritmo bajo dificultad controlada, aunque no representan el desempeño frente a datos reales~\parencite{Emmott2013, Goldstein2016}. Los escenarios guionados sobre el sistema real verifican la captura y el procesamiento de punta a punta. Los estudios con participantes que actúan en condiciones asignadas por consigna aportan datos de personas con etiquetas conocidas; en ellos, el diseño intrasujeto, en el que cada participante actúa como su propio control, reduce la variabilidad atribuible a diferencias individuales~\parencite{Charness2012}.

La factibilidad de un enfoque no supervisado descansa en que el modelo se ajusta sin ejemplos etiquetados de la clase anómala~\parencite{Chandola2009, Liu2008}: las etiquetas solo son necesarias para medir el desempeño, no para entrenar. Cuando los datos provienen de personas, su recolección queda sujeta al consentimiento informado y a los principios de protección de datos que se describen a continuación.

\subsection{Privacidad, protección de datos y minimización}
\label{sec:privacidad}

El tratamiento de datos conductuales en evaluaciones digitales requiere especial cuidado porque ciertas señales pueden utilizarse para verificar o confirmar la identidad de una persona. El Reglamento General de Protección de Datos de la Unión Europea define los datos biométricos como datos personales obtenidos mediante tratamiento técnico específico sobre características físicas, fisiológicas o conductuales que permiten o confirman la identificación única de una persona~\parencite{GDPR2016}.

En Argentina, la Ley N.\textsuperscript{o}~25.326 establece el régimen general de protección de datos personales y exige consentimiento libre, expreso e informado para el tratamiento de datos personales, salvo excepciones previstas por la norma~\parencite{Argentina2000Ley25326}. La Resolución AAIP N.\textsuperscript{o}~4/2019 incorporó criterios orientadores y definió datos biométricos como datos personales obtenidos a partir de tratamiento técnico específico sobre características físicas, fisiológicas o conductuales que permitan o confirmen la identificación única de una persona humana~\parencite{AAIP2019Res4}. La Ley N.\textsuperscript{o}~27.699 aprobó el Protocolo Modificatorio del Convenio 108 (Convenio 108+), que refuerza principios de finalidad, proporcionalidad, seguridad, transparencia, conservación limitada y derechos de los titulares de datos~\parencite{Argentina2022Ley27699}.

La normativa y la literatura asocian al tratamiento de datos conductuales en evaluaciones digitales los principios de privacidad desde el diseño: minimizar la captura, limitarla al entorno de evaluación, evitar datos ajenos a la consigna, separar identificadores directos de señales analíticas cuando sea posible y no utilizar resultados automatizados como base única para decisiones sancionatorias. Estas exigencias son consistentes con las advertencias de la literatura sobre \textit{online proctoring} y privacidad contextual~\parencite{Mutimukwe2023}.

\subsection{Síntesis del marco teórico}
\label{sec:sintesis_marco_teorico}

Los conceptos revisados delimitan un espacio intermedio entre la supervisión audiovisual intensiva y la evaluación digital que no ofrece información sobre el proceso de producción de las respuestas. La trazabilidad de autoría, apoyada en la biometría conductual de escritura y del uso del mouse y en señales contextuales de interacción, produce indicios del proceso que la literatura recomienda interpretar como evidencia complementaria, explicable y proporcional, sujeta a revisión humana y a los principios de protección de datos personales. El valor de esos indicios depende de las reglas de la evaluación y de la interpretación de quien la revisa.

Estos ejes sustentan la formulación de requerimientos y de arquitectura de los capítulos siguientes, en los que se los traduce en decisiones de diseño.
